\documentclass[10pt,a4paper]{amsart}
\usepackage[margin=19mm]{geometry}
\usepackage{amsmath,amssymb,mathtools}
\usepackage[scr=rsfs,cal=euler]{mathalfa}
\usepackage{xfrac}
\usepackage[unicode,hidelinks,bookmarks=false]{hyperref}
\newcommand{\ie}{i.e.\ }
\newcommand{\cf}{cf.\ }
\newcommand{\resp}{respectively\ }
\newcommand{\Subsection}{\subsection}
\providecommand{\nobreakdash}{\nobreak\hbox{-}}
\setlength{\emergencystretch}{2em}
\multlinegap0pt
\newtheorem{comment}{Comment}
\usepackage{amssymb}
\usepackage{float}
\usepackage[shortlabels]{enumitem}
\usepackage{siunitx}
\newtheorem{lemma}{Lemma}
\def\bPi{\mathbf{\Pi}}
\def\bkappa{\boldsymbol{\kappa}}
\def\bn{\boldsymbol{n}}
\newcommand\cF{\mathcal{F}}
\newcommand{\jump}[1]{[\![#1]\!]}
\def\rV{\mathrm{V}}
\newcommand\vdiv{\mathop{\mathrm{div}}\nolimits}
\title{A posteriori analysis of a virtual element approach on polytopal meshes for the buckling eigenvalue problem}
\author{Franco Dassi, Andres E Rubiano, Iv\'an Vel\'asquez}
\date{Source: arXiv:2603.20599v1 (CC BY 4.0)}
\begin{document}
\maketitle

\noindent\emph{Source and licence.} A posteriori analysis of a virtual element approach on polytopal meshes for the buckling eigenvalue problem. Version arXiv:2603.20599v1 (CC BY 4.0), distributed by arXiv under the Creative Commons Attribution 4.0 International licence, http://creativecommons.org/licenses/by/4.0/, as recorded in the arXiv licence metadata for this exact version and shown on the abstract page https://arxiv.org/abs/2603.20599v1. Full licence notice: \texttt{licenses/arxiv-2603.20599-CC-BY-4.0.txt}.\par\smallskip

\noindent\textit{Editorial scope.} Counted unit: the article's lemma ApVEPLem18 (source0.tex lines 847--852). The article's complete original proof of it is delivered with it (source0.tex lines 853--895, one \texttt{proof} environment of 43 lines). No mathematics is written, repaired or re-derived: the statement and the proof are the article's own lines, in the article's own notation, and no retained line is substituted or re-flowed.  Retained uncounted as the source-local context the proof needs: lines 257--259; lines 425--427; lines 651--660; lines 821--830; lines 833--840.  Nothing was omitted from any retained span, so the package carries no omission record.  No retained line contains a citation, so the wrapper carries no bibliography and no bibliography entry.  No retained line contains a figure environment, an image-inclusion command or drawing code, so no image is delivered and the package declares no asset dependency.  Editorial formatting only, in addition: an AMS \texttt{amsart} preamble that restates the article's own declarations and macros used by the retained lines, namely the environments \texttt{lemma} on separate counters, together with the standard packages those lines need.  The retained environments print in this document as Lemma 1 in order of appearance, whereas the article prints the same items with the numbering of its own full document.  The complete native source of the article as distributed in the arXiv e-print for this exact version is bundled unchanged as \texttt{sources/196862/source0.tex}.
\begin{align}\label{WMPr}
    a(u,v) = \lambda b(u,v), \quad \forall v\in \rV ,
\end{align}

\begin{align}\label{WDPr}
    a^h(u_h,v_h) = \lambda_h b^h(u_h,v_h), \quad \forall v_h\in \rV_h.
\end{align}

\begin{lemma}\label{Lem13ApVEP}
    If  $u$ and $u_h$ solve the spectral problems \eqref{WMPr} and \eqref{WDPr}, then for every $v\in  \rV$, the following identity holds:
    \begin{align*}
        a(u-u_h,v) &= \lambda b(u,v)-\lambda_h b( u_h,v) + \sum_{K\in \Omega^h} a_K(\Pi_K^{\nabla^2} u_h-u_h,v) \\
        &\quad + \lambda_h \sum_{K\in \Omega^h} \int_K (\bkappa(\nabla u_h - \bPi_K^0 (\nabla u_h)))\cdot \nabla v - \sum_{K\in \Omega^h} \int_{K}\lambda_h \vdiv\left( \bPi_K^0 (\bkappa \bPi_K^0(\nabla u_h))\right) v \\
        &\quad - \sum_{f\in \cF^h_\Omega} \int_f \jump{(\nabla^2\Pi_K^{\nabla^2} u_h )\bn_K^f}\cdot \nabla v + \sum_{f\in \mathcal{F}^h_\Omega} \int_f
        \jump{\lambda_h\bPi_K^0(\bkappa \bPi_K^0(\nabla u_h ))\cdot \bn_K^{f}}v \\
        &\quad + \sum_{K\in \Omega^h}\int_K \lambda_h \left(\bkappa \bPi_K^0(\nabla u_h)-\bPi_K^0(\bkappa \bPi_K^0(\nabla u_h))\right) \cdot \nabla v.
    \end{align*}
\end{lemma}

\begin{align}
||q_K||_{0,K}^2&\lesssim \int_K \psi_K q_K^2, &\quad \forall q_K\in \mathbb{P}_2(K),\label{bubb1}\\
||\psi_Kq_K||_{0,K} &\leq ||q_K||_{0,K}, &\quad \forall q_K\in \mathbb{P}_2(K),\label{bubb2}\\
h_f^{-1}||q_f||_{0,f}^2&\lesssim \int_f q_f \mathbf{n}^f_{K} \cdot \nabla(\psi_f q_f), &\quad \forall q_f\in \mathbb{P}_0(f),\label{bubb1Edge}\\
%
||\mathbf{n}^f_{K} \cdot \nabla(\psi_f q_f)||_{0,f} &\lesssim h^{-1}_f||q_f||_{0,f}, &\quad \forall q_f\in \mathbb{P}_0(f),\label{bubb2Edge}\\
\sum_{K=K^+,K^-} h_f^{-2}||\psi_f q_f||_{0,K}&\lesssim \sum_{K=K^+,K^-} |\psi_f q_f|_{2,K}, &\quad \forall q_f\in \mathbb{P}_0(f),\label{bubb3Edge} \\
\sum_{K=K^+,K^-} |\psi_f q_f|_{2,K}&\lesssim \sum_{K=K^+,K^-} h_f^{-2}||\psi_f q_f||_{0,K}, &\quad \forall q_f\in \mathbb{P}_0(f),\label{bubb3_5Edge} \\
\sum_{K=K^+,K^-} ||\psi_f q_f||_{0,K}&\lesssim h_f^{1/2}||q_f||_{0,f}, &\quad\forall q_f\in \mathbb{P}_0(f).\label{bubb4Edge}
\end{align}

\begin{align}
 | v|_{1,K}\lesssim h_K^{-1}||v||_{0,K} 
\quad
\mbox{and}
\quad 
 | v|_{2,K}\lesssim h_K^{-2}||v||_{0,K},
 &\quad \forall v\in  V_h^K .\label{bubb3}
\end{align}

\begin{lemma}\label{ApVEPLem18}
    If  $u$ and $u_h$ solve the spectral problems \eqref{WMPr} and \eqref{WDPr}, respectively,  then 
\begin{align}    
\Xi_K&\lesssim  |u-u_h|_{2,K} + S_K + h_K|\lambda u-\lambda_h u_h|_{1,K} + h_K\Lambda_K + h_K ||\nabla u_h - \bPi_K^0 (\nabla u_h)||_{0,K}.\nonumber
\end{align}
\end{lemma}

\begin{proof}
Let $K\in\Omega^h$. Consider $\psi_K$ as the interior bubble function defined in $K$ and  $v_K:=\psi_K\lambda_h\vdiv\left( \bPi_K^0 (\bkappa \bPi_K^0(\nabla u_h))\right)$. Since $v_K\in H_0^2(K)$, from now on we will denote its extension by $v$ by zero in $\Omega$. Thus, from the residual error equation in Lemma~\eqref{Lem13ApVEP} we readily see that
    \begin{align*} 
       \int_{K}\lambda_h \vdiv\left( \bPi_K^0 (\bkappa \bPi_K^0(\nabla u_h))\right) v  &= \lambda b_K(u,v)-\lambda_h b_K( u_h,v) + a_K(\Pi_K^{\nabla^2} u_h-u_h,v) -a_K(u-u_h,v)\\
        &\quad+ \lambda_h  \int_K (\bkappa(\nabla u_h - \bPi_K^0 (\nabla u_h)))\cdot \nabla v  \nonumber\\
        &\quad   + \int_K \lambda_h \left(\bkappa \bPi_K^0(\nabla u_h)-\bPi_K^0(\bkappa \bPi_K^0(\nabla u_h))\right) \cdot \nabla v
    \end{align*}
Next, we apply \eqref{bubb1} in the previous identity to obtain 
    \begin{align}
    ||\lambda_h \vdiv\left( \bPi_K^0 (\bkappa \bPi_K^0(\nabla u_h))\right) ||_{0,K}^2  &\lesssim \int_K \psi_K [\lambda_h \vdiv\left( \bPi_K^0 (\bkappa \bPi_K^0(\nabla u_h))\right)]^2 \nonumber \\ %= \int_{K}\lambda_h \vdiv\left( \bPi_K^0 (\bkappa \bPi_K^0(\nabla u_h))\right) v \nonumber\\
        &   = b(\lambda u- \lambda_h u_h,v) +  a_K(\Pi_K^{\nabla^2} u_h-u_h,v)-a(u-u_h,v) \nonumber\\
        &\quad + \lambda_h  \int_K (\bkappa(\nabla u_h - \bPi_K^0 (\nabla u_h)))\cdot \nabla v   \nonumber\\
        &\quad + \int_K \lambda_h \left(\bkappa \bPi_K^0(\nabla u_h)-\bPi_K^0(\bkappa \bPi_K^0(\nabla u_h))\right) \cdot \nabla v\nonumber\\
       & \lesssim |\lambda u-\lambda_h u_h|_{1,K}|v|_{1,K} + |\Pi_K^{\nabla^2} u_h-u_h|_{2,K} |v|_{2,K} \nonumber\\ 
       &\quad + |u-u_h|_{2,K}|v|_{2,K} + ||\nabla u_h - \bPi_K^0 (\nabla u_h)||_{0,K}|v|_{1,K}\nonumber\\
       &\quad + ||\lambda_h \left(\bkappa \bPi_K^0(\nabla u_h)-\bPi_K^0(\bkappa \bPi_K^0(\nabla u_h))\right) ||_{0,K}|v|_{1,K}.\label{cota_VOl_Eff}
    \end{align}
Since $v\in H_0^2(K)$, $0\leq \psi_K\leq 1$ and  $\lambda_h \vdiv\left( \bPi_K^0 (\bkappa \bPi_K^0(\nabla u_h))\right) \in \mathbb{P}_0(K)\in \rV_h$, the estimates \eqref{bubb2} and \eqref{bubb3} lead to 
\begin{subequations}
    \begin{align}
        |v|_{1,K} &\lesssim h_K^{-1}||\lambda_h \vdiv\left( \bPi_K^0 (\bkappa \bPi_K^0(\nabla u_h))\right) ||_{0,K}, \label{aux2-3_effa}\\ 
        |v|_{2,K} &\lesssim  h_K^{-2} ||\lambda_h \vdiv\left( \bPi_K^0 (\bkappa \bPi_K^0(\nabla u_h))\right) ||_{0,K}.\label{aux2-3_effb}
    \end{align}
\end{subequations}
Thus, by replacing the estimates \eqref{aux2-3_effa} and \eqref{aux2-3_effb} on the right-hand side of \eqref{cota_VOl_Eff}, and applying the definitions of $\Lambda_K$ and $S_K$ we have
\begin{align}    
||\lambda_h \vdiv\left( \bPi_K^0 (\bkappa \bPi_K^0(\nabla u_h))\right) ||_{0,K}
&\lesssim h_K^{-1}|\lambda u-\lambda_h u_h|_{1,K} + h_K^{-2}|\Pi_K^{\nabla^2} u_h-u_h|_{2,K} \nonumber\\
&\quad + h_K^{-2}|u-u_h|_{2,K} + h_K^{-1} ||\nabla u_h - \bPi_K^0 (\nabla u_h)||_{0,K} \nonumber\\
&\quad + h_K^{-1}||\lambda_h \left(\bkappa \bPi_K^0(\nabla u_h)-\bPi_K^0(\bkappa \bPi_K^0(\nabla u_h))\right) ||_{0,K} \nonumber\\
&\lesssim h_K^{-1}|\lambda u-\lambda_h u_h|_{1,K} + h_K^{-2} S_K + h_K^{-1} ||\nabla u_h - \bPi_K^0 (\nabla u_h)||_{0,K} \nonumber\\
&\quad + h_K^{-2}|u-u_h|_{2,K}  + h_K^{-1}\Lambda_K. \nonumber
\end{align}
%
Therefore, by using the definition of the local volume estimator $\Xi_K$ and multiplying by $h_K^{2}$ in the above inequality, we obtain the result.
%
\begin{comment}
\begin{align}    
\Xi_K%=&h_K^{2}||\lambda_h \vdiv\left( \bPi_K^0 (\bkappa \bPi_K^0(\nabla u_h))\right) ||_{0,K}\nonumber\\
&\lesssim h_K^{2}|\lambda u-\lambda_h u_h|_{1,K} +  S_K + h_K ||\nabla u_h - \bPi_K^0 (\nabla u_h)||_{0,K}+ |u-u_h|_{2,K}  + h_K\Lambda_K.\nonumber
\end{align}
\end{comment}
\end{proof}

\end{document}
